\documentclass[11pt]{article}
\usepackage[margin=1in]{geometry}
\usepackage{amsmath}
\title{Orbits and energy}
\author{Physics 201, week 6}
\date{}
\begin{document}
\maketitle
\section{Kepler's third law}
For a planet of mass $m$ circling a star of mass $M \gg m$ at radius $r$, gravity supplies the centripetal force:
\begin{equation}
  \frac{G M m}{r^2} = \frac{m v^2}{r}, \qquad v = \frac{2\pi r}{T}.
\end{equation}
Solving for the period gives
\begin{equation}
  T^2 = \frac{4\pi^2}{G M}\, r^3 .
\end{equation}
\section{Energy of an orbit}
The total energy is the sum of kinetic and potential energy,
\begin{equation}
  E = \tfrac12 m v^2 - \frac{G M m}{r} = -\frac{G M m}{2r},
\end{equation}
so a lower orbit is more tightly bound.
\end{document}
